\section*{Chapter \ref{ch:b2:planet-formation-meteorites} --- Planet Formation and Meteorites}

\begin{solution}{exo:b2:planet-formation-meteorites:1}
Gas: \qty{2200}{g.cm^{-2}} at \qty{1}{au}; at \qty{5.2}{au},
$2200\times5.2^{-3/2} = 2200/11.86 = \qty{186}{g.cm^{-2}}$. Solids:
$33\times0.33 = \qty{10.9}{g.cm^{-2}}$ at \qty{1}{au};
$33/11.86 = \qty{2.78}{g.cm^{-2}}$ at \qty{5.2}{au}. At \qty{5.2}{au} the gas is
$186/2.78 = 67$ times heavier than the solids.
\end{solution}

\begin{solution}{exo:b2:planet-formation-meteorites:2}
Aluminium (\qty{1653}{K}, ultrarefractory), magnesium (\qty{1336}{K}, common),
iron (\qty{1334}{K}, common), sodium (\qty{958}{K}, moderately volatile), sulfur
(\qty{664}{K}, volatile), oxygen as water ice (\qty{180}{K}, highly volatile).
CAIs, the first condensates, are enriched in aluminium (and calcium, titanium);
the Earth is depleted in the volatile ones, sodium, sulfur and water, relative to
magnesium and aluminium.
\end{solution}

\begin{solution}{exo:b2:planet-formation-meteorites:3}
Falls: ordinary \omterm{def:b2:planet-formation-meteorites:chondrule}{chondrites} $856/1107 = \qty{77.3}{\percent}$, irons
$49/1107 = \qty{4.4}{\percent}$. Finds: $38647/44534 = \qty{86.8}{\percent}$ and
$1020/44534 = \qty{2.3}{\percent}$. Antarctic finds: $22096/44534 =
\qty{49.6}{\percent}$. The falls: every observed fall is collected whatever its
class, whereas a find depends on how well the \omterm{def:b2:planet-formation-meteorites:fall-find}{meteorite} survives on the ground,
how it looks and where people search, and half of the finds come from one kind of
place.
\end{solution}

\begin{solution}{exo:b2:planet-formation-meteorites:4}
$\Delta t = 1.035\ln(\num{5.23e-5}/\num{1.0e-5}) = 1.035\ln 5.23 =
1.035\times1.654 = \qty{1.71}{Myr}$ after the CAIs.
\end{solution}

\begin{solution}{exo:b2:planet-formation-meteorites:5}
$4\pi\Sigma_1 r_1^2 = 4\pi\times\num{2.2e4}\times(\num{1.496e11})^2 =
\qty{6.19e27}{kg}$. Inside \qty{30}{au}: $\times\sqrt{30} = 5.48$, so
\qty{3.39e28}{kg} $= 0.017$ solar masses; inside \qty{100}{au}: $\times10$,
\qty{6.19e28}{kg} $= 0.031$ solar masses. Both lie within the observed range,
nearer its upper end: the nebula was an ordinary disc.
\end{solution}

\begin{solution}{exo:b2:planet-formation-meteorites:6}
$r = (120/T)^{7/3}$ au: for water, $(120/182)^{7/3} = 0.659^{7/3} =
\qty{0.38}{au}$; for methane, $(120/41)^{7/3} = 2.93^{7/3} = \qty{12.3}{au}$.
The grain in direct sunlight is the surface layer of the disc; the midplane,
shielded by the dust above it, is colder, so ice survives closer in. The
\omterm{def:b2:planet-formation-meteorites:planetesimal}{planetesimals} formed in the midplane, so the midplane estimate applies, provided
the disc is passive; heating by accretion moves the \omterm{def:b2:planet-formation-meteorites:snow-line}{snow line} out.
\end{solution}

\begin{solution}{exo:b2:planet-formation-meteorites:7}
At \qty{1}{au}: $r^2 = \num{2.238e22}$, $\Sigma_p = \qty{109}{kg.m^{-2}}$, so
$2\pi b r^2\Sigma_p = 31.42\times\num{2.238e22}\times109 = \num{7.66e25}$;
$(\num{7.66e25})^{3/2} = \num{6.71e38}$, divided by
$(3\times\num{1.989e30})^{1/2} = \num{2.443e15}$: \qty{2.75e23}{kg}
$= 0.046$ Earth masses. At \qty{5.2}{au}: $r^2 = \num{6.05e23}$,
$\Sigma_p = \qty{27.8}{kg.m^{-2}}$, $2\pi b r^2\Sigma_p = \num{5.29e26}$,
$M = \num{1.22e40}/\num{2.443e15} = \qty{4.98e24}{kg} = 0.83$ Earth masses.
With \num{3.5} times more solids, $M$ grows as $3.5^{3/2} = 6.55$: 5.5 Earth
masses.
\end{solution}

\begin{solution}{exo:b2:planet-formation-meteorites:8}
Per kilogram, $\num{1.01e19}$ atoms, so
$\num{1.01e19}\times\num{5.00e-13} = \qty{5.05e6}{J.kg^{-1}}$. The temperature
rise would be $\num{5.05e6}/837 = \qty{6030}{K}$. Warming from \num{250} to
\qty{1700}{K} takes $837\times1450 = \qty{1.21e6}{J.kg^{-1}}$: aluminium-26
supplied about four times that.
\end{solution}

\begin{solution}{exo:b2:planet-formation-meteorites:9}
$\num{5.07e10}/\num{7.49e7} = 677$. The mean life of iron-60 is
$2.62/\ln2 = \qty{3.78}{Myr}$. Initial powers per gram:
$\num{5.07e10}/1.035 = \num{4.90e10}$ and $\num{7.49e7}/3.78 = \num{1.98e7}$ erg per
gram per million years, a ratio of about \num{2470}: aluminium-26 dominated by far.
\end{solution}

\begin{solution}{exo:b2:planet-formation-meteorites:10}
The powers are equal when $\num{4.90e10}\,e^{-t/1.035} = \num{1.98e7}\,e^{-t/3.78}$,
that is $t(1/1.035 - 1/3.78) = \ln(\num{4.90e10}/\num{1.98e7}) = \ln 2475 = 7.81$;
with $0.966 - 0.265 = 0.702$ per million years, $t = \qty{11.1}{Myr}$. By then
iron-60 has left $\num{7.49e7}\,e^{-11.1/3.78} = \num{7.49e7}\times0.053 =
\qty{4.0e6}{erg.g^{-1}}$, that is \qty{400}{J.kg^{-1}}: thousands of times too
little to melt anything.
\end{solution}

\begin{solution}{exo:b2:planet-formation-meteorites:11}
$\tau = 1.035\times\num{3.156e13} = \qty{3.267e13}{s}$ and
$H_0\tau = \qty{6.21e6}{J.kg^{-1}}$. For \qty{1200}{K}:
$t_a^* = 1.035\ln\bigl(\num{6.21e6}/(837\times950)\bigr) = 1.035\ln7.81 =
\qty{2.13}{Myr}$; for \qty{1700}{K}: $1.035\ln\bigl(\num{6.21e6}/(837\times1450)\bigr)
= 1.035\ln5.12 = \qty{1.69}{Myr}$. Halving $H_0$ subtracts $\tau\ln2 =
\qty{0.72}{Myr}$ from both, since $t_a^* = \tau\ln H_0 + \tau\ln(\tau/c(T_c -
T_0))$ and only the first term changes: the shift is one half-life, whatever the
critical temperature.
\end{solution}

\begin{solution}{exo:b2:planet-formation-meteorites:12}
$\kappa = 2.1/(3300\times837) = \qty{7.60e-7}{m^2.s^{-1}}$ and
$\sqrt{\kappa\tau} = \sqrt{\num{7.60e-7}\times\num{3.267e13}} = \qty{4980}{m}$,
about \qty{5}{km}. The centre of a body \qty{100}{km} across is \qty{50}{km} deep,
ten times this distance: its heat cannot escape while aluminium-26 decays, and the
proposition applies. A body \qty{2}{km} across loses its heat as fast as it makes
it. Melting needs a radius of a few times \qty{5}{km}, that is a diameter of some
\qtyrange{10}{20}{km} or more, as the models find.
\end{solution}

\begin{solution}{pb:b2:planet-formation-meteorites:1}
\textbf{1.} A fall: its arrival was seen and it was picked up the next morning.
It is unweathered, and it adds to the unbiased sample of what reaches the Earth.

\textbf{2.} $856/1107 = 0.77$: about three chances in four.

\textbf{3.} \omterm{def:b2:planet-formation-meteorites:chondrule}{Chondrules} \qty{0.3}{mm} across, \qty{9}{\percent} metal and olivine with
\qty{18}{\percent} fayalite: an H \omterm{def:b2:planet-formation-meteorites:chondrule}{chondrite} (\qty{10}{\percent} metal, Fa 16--20).

\textbf{4.} Blurred \omterm{def:b2:planet-formation-meteorites:chondrule}{chondrules} and uniform minerals mean thermal metamorphism: types
5 to 6. The ordinary \omterm{def:b2:planet-formation-meteorites:chondrule}{chondrites} reached at most about \qty{1220}{K}, that is
\qty{947}{\degreeCelsius}.

\textbf{5.} Ordinary \omterm{def:b2:planet-formation-meteorites:chondrule}{chondrites} hold only \qtyrange{0.1}{1}{\percent} CAIs by
volume. The stone is made of \omterm{def:b2:planet-formation-meteorites:chondrule}{chondrules} and matrix assembled after the CAIs, which
are the oldest solids: any rock that contains them is younger.

\textbf{6.} Each mineral received at formation the same magnesium plus aluminium-26
in proportion to its aluminium; the aluminium-26 has since turned into
magnesium-26, so the more aluminium a mineral had per magnesium, the more
magnesium-26 it gained.

\textbf{7.} $(0.139590 - 0.139500)/(15.0 - 0.05) = \num{0.000090}/14.95 =
\num{6.02e-6}$: the ratio $\ce{^{26}Al}/\ce{^{27}Al}$ of the \omterm{def:b2:planet-formation-meteorites:chondrule}{chondrule} when it
formed.

\textbf{8.} $0.139500 + \num{6.02e-6}\times3.00 = 0.139518$: on the line.

\textbf{9.} $\Delta t = 1.035\ln(\num{5.23e-5}/\num{6.02e-6}) = 1.035\ln8.69 =
1.035\times2.162 = \qty{2.24}{Myr}$.

\textbf{10.} $4567.5 - 2.2 = \qty{4565.3}{\Ma}$.

\textbf{11.} $\num{5.23e-5}\,e^{-10/1.035} = \num{5.23e-5}\times\num{6.4e-5} =
\num{3.3e-9}$, far too little magnesium-26 to measure. Hafnium-182 (half-life
\qty{8.9}{Myr}) or manganese-53 (\qty{3.7}{Myr}) suit events some ten million years
after the CAIs.

\textbf{12.} $\kappa = 2.1/(3300\times837) = \qty{7.60e-7}{m^2.s^{-1}}$;
$R^2/\kappa = \num{4e8}/\num{7.60e-7} = \qty{5.26e14}{s} = \qty{16.7}{Myr}$, sixteen
times the mean life of aluminium-26: the centre keeps its heat.

\textbf{13.} $\num{1.9e-7}\,e^{-2.4/1.035} = \num{1.9e-7}\times0.0984 =
\qty{1.87e-8}{W.kg^{-1}}$.

\textbf{14.} $\num{1.87e-8}\times\num{3.267e13} = \qty{6.11e5}{J.kg^{-1}}$.

\textbf{15.} $\Delta T = \num{6.11e5}/837 = \qty{730}{K}$; peak
$250 + 730 = \qty{980}{K}$.

\textbf{16.} Above \qty{750}{K}, so the rock sintered and was metamorphosed; below
\qty{1200}{K}, so neither metal nor sulfide melted. The stone is a heated \omterm{def:b2:planet-formation-meteorites:chondrule}{chondrite},
of high \omterm{def:b2:planet-formation-meteorites:type}{petrologic type}, below the \qty{1220}{K} of question 4: consistent.

\textbf{17.} $1.035\ln\bigl(\num{6.21e6}/(837\times950)\bigr) = \qty{2.13}{Myr}$:
about \qty{0.3}{Myr} earlier.

\textbf{18.} The rise is $\Delta T(1 - e^{-(t - t_a)/\tau})$; \qty{90}{\percent}
needs $e^{-(t - t_a)/\tau} = 0.1$, $t - t_a = 1.035\ln10 = \qty{2.38}{Myr}$, that is
\qty{4.8}{Myr} after the CAIs.

\textbf{19.} $\Delta T = 7420\,e^{-0.8/1.035} = 7420\times0.462 = \qty{3430}{K}$,
peak \qty{3680}{K}, far above the \qty{1700}{K} at which the silicates are molten.
The body melted (the heat of melting held its temperature down) and
differentiated into a metal core and a silicate mantle and crust.

\textbf{20.} $e^{(2.4 - 0.8)/1.035} = e^{1.546} = 4.7$ times more.

\textbf{21.} Zinc and lead condense near \qty{727}{K}, moderately volatile; magnesium
(\qty{1336}{K}) and aluminium (\qty{1653}{K}) condense far earlier. The solids of
the stone's region gathered before the gas had cooled enough for all the zinc and
lead to condense.

\textbf{22.} $33\times0.33\times2.5^{-3/2} = 10.89/3.95 = \qty{2.76}{g.cm^{-2}}$,
\qty{27.6}{kg.m^{-2}}.

\textbf{23.} $r^2 = (3.74\times10^{11})^2 = \num{1.40e23}$;
$2\pi\times5\times\num{1.40e23}\times27.6 = \num{1.21e26}$;
$(\num{1.21e26})^{3/2}/\num{2.443e15} = \qty{5.45e23}{kg}$, $0.091$ Earth masses.
The parent body, $\tfrac43\pi(\num{2e4})^3\times3300 = \qty{1.1e17}{kg}$, is five
million times lighter: a \omterm{def:b2:planet-formation-meteorites:planetesimal}{planetesimal}, not an embryo.

\textbf{24.} The centre of the parent body peaked at about
\textbf{\qty[text-series-to-math=true]{707}{\degreeCelsius}} ($980 - 273$): hot enough to
recrystallise the \omterm{def:b2:planet-formation-meteorites:chondrule}{chondrite}, too cool to melt its metal.
\end{solution}
